 \documentclass[10pt]{beamer}

% ===============================
% Theme Configuration
% ===============================
\usetheme{Madrid}
\usecolortheme{default}
\setbeamertemplate{navigation symbols}{}     % Remove navigation icons
\setbeamertemplate{footline}[frame number]   % Slide numbers in footer

% ===============================
% Packages
% ===============================
\usepackage{graphicx}
\usepackage{booktabs}
\usepackage{amsmath}
\usepackage{array}
\usepackage{tikz}
\usepackage{hyperref}

% ===============================
% Title Page Information
% ===============================
\title[Crowd Presence & Identity Detection]
{Crowd Presence and Identity Detection and Recognition}

\subtitle{B.Tech Mini Project}

\author{
ALEN ALIAS (FIT23CS025) \\
ASHLIN P J (FIT23CS049) \\
AV NAVNEETH SAGAR (FIT23CS055) \\
BENSON ELDHO (FIT23CS059)
}

\institute{
Department of Computer Science and Engineering \\
Federal Institute of Science and Technology \\[6pt]
\textbf{Project Guide:} Nobel G
}

\date{}   % Date removed

% ===============================
% College Logo
% ===============================
\titlegraphic{
\includegraphics[width=2cm]{logo.png}
}

% ===============================
% Document Start
% ===============================
\begin{document}

% ===============================
% Title Page Frame
% ===============================
\begin{frame}
\titlepage
\end{frame}
\begin{frame}{Introduction}

\begin{itemize}
    \item \textbf{Crowd Presence and Identity Detection \& Recognition} is a real-time computer vision system designed to monitor and verify human presence in controlled environments.
    
    \item It is implemented in contexts such as \textbf{college canteens} to ensure accurate tracking of user visits and activities.
    
    \item The system captures input from \textbf{live camera feeds or RTSP streams} and processes it continuously.
    
    \item It performs \textbf{face detection and identity recognition} by comparing captured faces with registered user data.
    
    \item All recognized events are stored as \textbf{timestamped records}, ensuring traceability and accountability.
\end{itemize}

\vspace{0.3cm}

\textbf{Core Components:}
\begin{itemize}
    \item Detection: YOLOv8 / OpenCV DNN
    \item Recognition: DeepFace (FaceNet512 embeddings)
    \item Storage: SQLite Database
\end{itemize}

\vspace{0.2cm}

\begin{itemize}
    \item The system enhances \textbf{transparency} and helps prevent disputes by maintaining reliable digital logs.
\end{itemize}

\end{frame}
\begin{frame}{Problem Statement}

\begin{itemize}
    \item Lack of a reliable system to \textbf{verify actual presence} of individuals in monitored environments.
    
    \item Possibility of \textbf{false claims or disputes} due to absence of verifiable evidence.
    
    \item Existing methods rely on \textbf{manual record keeping}, which is time-consuming and error-prone.
    
    \item No automated mechanism to \textbf{track and log user activity} in real time.
    
    \item Limited access to \textbf{analytics} such as peak hours, usage patterns, and visitor trends.
\end{itemize}

\vspace{0.3cm}

\textbf{Need:}
\begin{itemize}
    \item An \textbf{automated, intelligent, and scalable system} for real-time presence detection and identity recognition.
    
    \item A solution that provides \textbf{accurate logging, transparency, and data-driven insights}.
\end{itemize}

\end{frame}
\begin{frame}{Project Objectives}

\begin{itemize}
    \item Develop a \textbf{real-time face detection and recognition system} for monitoring human presence.
    
    \item Accurately \textbf{identify registered individuals} using deep learning-based face embeddings.
    
    \item Automatically \textbf{log presence events} with timestamps and relevant metadata.
    
    \item Detect and handle \textbf{unknown or unregistered individuals}.
    
    \item Provide \textbf{statistical insights and reports} such as usage patterns and peak activity periods.
    
    \item Support multiple input sources including \textbf{webcam and RTSP streams}.
    
    \item Ensure a \textbf{secure, efficient, and scalable architecture} for real-world deployment.
\end{itemize}

\end{frame}
\begin{frame}{Purpose and Need of the Project}

\textbf{Purpose:}
\begin{itemize}
    \item Maintain \textbf{transparent and verifiable records}.
    \item Provide \textbf{evidence for dispute resolution}.
    \item Improve \textbf{monitoring efficiency}.
\end{itemize}

\vspace{0.3cm}

\textbf{Need:}
\begin{itemize}
    \item Replace \textbf{manual record systems}.
    \item Enable \textbf{automated real-time tracking}.
    \item Support \textbf{data-driven decision making}.
\end{itemize}

\end{frame}
\begin{frame}{Scope of the Project}

\textbf{Includes:}
\begin{itemize}
    \item Real-time \textbf{face detection and recognition} using webcam or RTSP streams.
    
    \item \textbf{User registration} with face capture from images or live camera input.
    
    \item Automated \textbf{presence/visit logging} with timestamps and duration tracking.
    
    \item \textbf{GUI-based interface} for monitoring, management, and interaction.
    
    \item Generation of \textbf{reports and CSV exports} for analysis and record keeping.
\end{itemize}

\vspace{0.3cm}

\textbf{Excludes:}
\begin{itemize}
    \item \textbf{Cloud-based deployment} and remote server integration.
    
    \item Integration with other \textbf{biometric systems} (e.g., fingerprint, iris).
    
\end{itemize}

\end{frame}
\begin{frame}{Social Relevance and SDGs}

\textbf{Social Impact:}
\begin{itemize}
    \item Promotes \textbf{smart and digital monitoring} in institutional environments.
    
    \item Enhances \textbf{transparency and accountability} through verifiable records.
    
    \item Improves \textbf{safety and security} by tracking presence of individuals.
    
    \item Supports efficient \textbf{resource and facility management}.
\end{itemize}

\vspace{0.3cm}

\textbf{Aligned with UN Sustainable Development Goals:}
\begin{itemize}
    \item \textbf{SDG 4 – Quality Education:} Enables better campus management systems.
    
    \item \textbf{SDG 9 – Industry, Innovation and Infrastructure:} Uses AI-based smart solutions.
    
    \item \textbf{SDG 16 – Peace, Justice and Strong Institutions:} Ensures transparency and reliable records.
\end{itemize}

\end{frame}
\begin{frame}{Software Requirement Specification}

\begin{itemize}
    \item Defines the \textbf{requirements and expectations} of the system for proper design and implementation.
    
    \item Covers both \textbf{functional} and \textbf{non-functional} aspects to ensure system completeness.
    
    \item Specifies how the system interacts with \textbf{users, hardware, and data sources}.
    
    \item Acts as a \textbf{foundation for development, testing, and validation}.
    
    \item Ensures the system meets goals of \textbf{accuracy, reliability, and performance}.
\end{itemize}

\vspace{0.3cm}

\textbf{Categories:}
\begin{itemize}
    \item Functional Requirements (System behavior and features)
    \item Non-Functional Requirements (Performance and quality constraints)
\end{itemize}

\end{frame}
\begin{frame}{Functional Requirements}

\begin{itemize}
    \item Capture video input from \textbf{webcam or RTSP stream} for continuous monitoring.
    
    \item Perform \textbf{real-time face detection} on incoming frames.
    
    \item Recognize \textbf{registered individuals} using stored face embeddings.
    
    \item Support \textbf{user registration} through image upload or live capture.
    
    \item Automatically \textbf{log presence events} with timestamps and identity details.
    
    \item Detect and store data of \textbf{unknown individuals} for further analysis.
    
    \item Provide \textbf{GUI/CLI interfaces} for system control and monitoring.
    
    \item Generate \textbf{reports and export data} (e.g., CSV) for analysis.
\end{itemize}

\end{frame}
\begin{frame}{Non-Functional Requirements}

\begin{itemize}
    \item Ensure \textbf{low response time} for near real-time detection and recognition.
    
    \item Maintain \textbf{high accuracy} in face recognition under varying lighting and angles.
    
    \item Provide \textbf{reliability and stability} during continuous operation.
    
    \item Ensure \textbf{data security and privacy} for stored images and embeddings.
    
    \item Support \textbf{scalability} to handle increasing number of users and records.
    
    \item Maintain \textbf{usability} through an intuitive GUI and simple CLI interface.
    
    \item Optimize \textbf{performance} for both CPU and GPU environments.
\end{itemize}

\end{frame}
\begin{frame}{System Architecture}

\begin{center}
\includegraphics[width=0.53\linewidth]{architecture.png}
\end{center}

\vspace{0.2cm}

\textbf{Flow:} Camera/RTSP (go2rtc relay) → Detection (OpenCV DNN) → DeepFace (SFace embeddings) → Matching \& Visit Logger → SQLite/GUI

\end{frame}
\begin{frame}{Application Architecture Design}
\begin{center}
    \includegraphics[width=0.54\linewidth]{Application_architecture.png}
\end{center}

\end{frame}
\begin{frame}{Application Architecture Design}

\begin{itemize}
    \item The system follows a \textbf{modular, layered architecture} with clear separation of responsibilities.
    \item Organized into key runtime modules: \textbf{GUI, Camera Pipeline, Recognition Pipeline, and Database Layer}.
    \item Recognition and capture run with \textbf{threaded processing} to keep UI responsive.
\end{itemize}

\vspace{0.3cm}

\textbf{Core Components:}
\begin{itemize}
    \item \textbf{UI Layer:} Tkinter GUI and CLI for monitoring, registration, and log export.
    \item \textbf{Camera Pipeline:} Single/multi-camera capture, RTSP reconnect, low-latency buffering.
    \item \textbf{Recognition Layer:} Detection, embedding extraction, gallery-based matching, stable labels.
    \item \textbf{Data Layer:} SQLite tables for users, visit logs, unknown faces, and statistics.
\end{itemize}

\vspace{0.3cm}

\begin{itemize}
    \item Ensures \textbf{scalability, maintainability, and easy extension} of system features.
\end{itemize}

\end{frame}
\begin{frame}{GUI Design}

\begin{center}
\includegraphics[width=0.8\linewidth]{gui_design.png}
\end{center}

\vspace{0.2cm}

Modules:
\begin{itemize}
    \item Live Detection
    \item Student Management
    \item Visit Logs
    \item Statistics Dashboard
\end{itemize}

\end{frame}
\begin{frame}{API Design}

\begin{itemize}
    \item The system uses \textbf{internal APIs (service interfaces)} for communication between modules.
    
    \item APIs enable interaction between \textbf{UI, Services, and Data layers}.
    
    \item Designed to ensure \textbf{modularity, reusability, and maintainability}.
\end{itemize}

\vspace{0.3cm}

\textbf{Key API Operations:}
\begin{itemize}
    \item \textbf{Detection API:} Process frames and detect faces.
    
    \item \textbf{Recognition API:} Match detected faces with stored embeddings.
    
    \item \textbf{Registration API:} Add/update user data and face information.
    
    \item \textbf{Visit Logging API:} Record presence events with timestamps.
    
    \item \textbf{Report API:} Generate logs and export data (CSV).
\end{itemize}

\end{frame}
\begin{frame}{Database Design}

\begin{itemize}
    \item The system uses a \textbf{SQLite database} for efficient data management.
    
    \item It stores \textbf{user information, visit logs, and recognition data}.
    
    \item Maintains a \textbf{1:N relationship} between users and visit records.
\end{itemize}
\begin{center}
    \includegraphics[width=0.5\linewidth]{Database design diagram overview.png}
\end{center}

\end{frame}
\begin{frame}{Technology Stack}

\begin{center}
\begin{tabular}{|l|l|}
\hline
\textbf{Component} & \textbf{Technology Used} \\ \hline
Programming Language & Python 3.10+ \\ \hline
Face Detection & YOLOv8 / OpenCV DNN \\ \hline
Face Recognition & DeepFace (SFace embeddings) \\ \hline
Image Processing & OpenCV \\ \hline
AI Backend & TensorFlow \\ \hline
Database & SQLite \\ \hline
GUI & Tkinter \\ \hline
RTSP Relay & go2rtc \\ \hline
\end{tabular}
\end{center}

\end{frame}
\begin{frame}{System Modules Overview}

\begin{itemize}
    \item The system is divided into \textbf{independent modules} for better organization.
    
    \item Each module handles a specific \textbf{functional responsibility}.
    
    \item Modules interact through \textbf{service-based architecture}.
\end{itemize}

\vspace{0.3cm}

\textbf{Main Modules:}
\begin{itemize}
    \item Camera \& Stream Module (webcam/RTSP/go2rtc)
    \item Detection \& Recognition Module
    \item Registration Module (live, image, multi-photo, bulk CSV)
    \item Visit Logging \& Reporting Module
\end{itemize}

\end{frame}
\begin{frame}{Detection Module}
\begin{center}
    \includegraphics[width=0.8\linewidth]{Detection module.png}
\end{center}

\end{frame}
\begin{frame}{Recognition Module}

\begin{center}
    \includegraphics[width=0.75\linewidth]{Recognition module.png}
\end{center}

\end{frame}
\begin{frame}{Registration Module}

\begin{center}
    \includegraphics[width=0.75\linewidth]{Registration module.png}
\end{center}

\end{frame}
\begin{frame}{Logging \& Reporting Module}

\begin{center}
    \includegraphics[width=0.75\linewidth]{Logging and reporting.png}
\end{center}

\end{frame}

\begin{frame}{Face Detection Algorithm}

\begin{itemize}
    \item Capture frame from \textbf{webcam or RTSP stream}.
    
    \item Preprocess frame (resize, normalization).
    
    \item Apply \textbf{OpenCV DNN} (with fallback detection) for face localization.
    
    \item Filter detections based on \textbf{confidence threshold}.
    
    \item Extract and validate face regions for recognition queue.
\end{itemize}

\end{frame}
\begin{frame}{Face Recognition Algorithm}

\begin{itemize}
    \item Receive detected face from detection module.
    
    \item Generate \textbf{face embedding} using DeepFace (SFace model).
    \item Compare against \textbf{gallery embeddings} of each registered user.
    \item Compute cosine similarity and apply threshold + margin checks.
    \item Stabilize labels across frames before confirming \textbf{known/unknown}.
\end{itemize}

\end{frame}
\begin{frame}{User Registration Procedure}

\begin{itemize}
    \item Capture samples from \textbf{live camera}, \textbf{single image}, or \textbf{multi-photo set}.
    
    \item Detect and extract valid face regions.
    
    \item Generate embeddings for each valid sample.
    
    \item Store \textbf{multi-sample gallery embeddings} and compute average embedding.
    
    \item Support \textbf{bulk CSV registration} with per-user photo folders.
\end{itemize}

\end{frame}
\begin{frame}{Visit Logging Procedure}

\begin{itemize}
    \item Receive stabilized recognized identity from recognition module.
    
    \item Apply \textbf{cooldown period} to avoid duplicate rapid logs.
    
    \item Log entry time and optionally capture screenshot path.
    
    \item Store known/unknown status for statistics and filtering.
    \item Persist record in \textbf{visit\_logs table}.
\end{itemize}

\end{frame}
\begin{frame}{Frame Processing Pipeline}

\begin{itemize}
    \item Continuously fetch frames using dedicated \textbf{capture threads}.
    \item Keep low-latency latest-frame buffering for responsive UI.
    \item Perform detection and recognition with frame-skip optimizations.
    \item Push crops to background \textbf{recognition queue} and merge results asynchronously.
    \item Render annotated results and live counters on GUI.
\end{itemize}

\end{frame}
\begin{frame}{Report Generation Procedure}

\begin{itemize}
    \item Fetch filtered data from \textbf{database}.
    
    \item Process logs based on date/user filters.
    
    \item Compute statistics (visits, duration).
    
    \item Generate structured report.
    
    \item Export data in \textbf{CSV format}.
\end{itemize}
\end{frame}
\begin{frame}{Coding Standards, Environment Setup \& Source Code Control}

\textbf{Coding Standards:}
\begin{itemize}
    \item Follow \textbf{PEP8 guidelines} for Python code.
    \item Use meaningful \textbf{variable and function names}.
    \item Modularize code into \textbf{services, repositories, and utilities}.
    \item Include \textbf{comments and docstrings} for clarity.
\end{itemize}

\vspace{0.3cm}

\textbf{Environment Setup:}
\begin{itemize}
    \item Python 3.10+ installed.
    \item Install dependencies using \textbf{requirements.txt}.
    \item Set up SQLite database and required directories.
    \item GPU support configured via TensorFlow (if available).
\end{itemize}

\vspace{0.3cm}

\textbf{Source Code Control Setup:}
\begin{itemize}
    \item Use \textbf{Git} for version control.
    \item Maintain separate \textbf{branches} for features and bug fixes.
    \item Regular commits with descriptive \textbf{commit messages}.
    \item Backup code on \textbf{GitHub or institutional server}.
\end{itemize}

\end{frame}

% Slide 27: Unit Testing
\begin{frame}{Unit Testing (All Core Functions)}
\textbf{Objective:} Validate each module independently.

\vspace{0.2cm}
\footnotesize
\begin{center}
\begin{tabular}{|p{3.8cm}|p{2.1cm}|p{3.0cm}|}
\hline
\textbf{Unit Area} & \textbf{Cases} & \textbf{Result} \\ \hline
Face Matching Functions & 8 & 8/8 Passed \\ \hline
Database Functions & 10 & 10/10 Passed \\ \hline
Camera Pipeline Functions & 6 & 6/6 Passed \\ \hline
Utility Functions & 7 & 7/7 Passed \\ \hline
\textbf{Total} & \textbf{31} & \textbf{31/31 Passed} \\ \hline
\end{tabular}
\end{center}
\normalsize

\textbf{Tools:} \texttt{unittest}, \texttt{test\_system.py}, \texttt{test\_rtsp.py}
\end{frame}

% Slide 28: Integration Testing
\begin{frame}{Integration Testing}
\textbf{Objective:} Verify module-to-module data flow.

\vspace{0.2cm}
\footnotesize
\begin{center}
\begin{tabular}{|p{3.7cm}|p{3.8cm}|p{1.5cm}|}
\hline
\textbf{Integration Path} & \textbf{Expected Behavior} & \textbf{Status} \\ \hline
Camera $\rightarrow$ Detection & face boxes generated correctly & Pass \\ \hline
Detection $\rightarrow$ Recognition & known/unknown label assignment & Pass \\ \hline
Recognition $\rightarrow$ Logger & cooldown + timestamp logging & Pass \\ \hline
GUI $\leftrightarrow$ Database & registration, logs, CSV export & Pass \\ \hline
go2rtc $\rightarrow$ Multi-camera UI & multiple relay stream ingest and rendering & Pass \\ \hline
\end{tabular}
\end{center}
\normalsize
\end{frame}

% Slide 29: Functional Testing
\begin{frame}{Functional Testing}
\textbf{Objective:} Validate end-user features from GUI/CLI.

\vspace{0.2cm}
\footnotesize
\begin{center}
\begin{tabular}{|p{3.6cm}|p{3.9cm}|p{1.6cm}|}
\hline
\textbf{Feature} & \textbf{Validation} & \textbf{Result} \\ \hline
Live Monitoring & identity overlay on active feed & OK \\ \hline
Multiple Camera Mode & simultaneous multi-stream display & OK \\ \hline
User Registration & live/image/bulk CSV onboarding & OK \\ \hline
Visit Logs \& Export & filter + CSV report generation & OK \\ \hline
Dashboard & counts and trends refresh & OK \\ \hline
\end{tabular}
\end{center}
\normalsize
\end{frame}

% Slide 30: Load Testing
\begin{frame}{Load Testing}
\textbf{Objective:} Check stability under continuous and concurrent load.

\vspace{0.2cm}
\footnotesize
\begin{center}
\begin{tabular}{|p{3.1cm}|p{3.0cm}|p{2.9cm}|}
\hline
\textbf{Scenario} & \textbf{Load} & \textbf{Outcome} \\ \hline
Single Stream Run & 30 min continuous RTSP & Stable, no freeze \\ \hline
Multi-stream Run & multiple camera relay streams via go2rtc & Responsive UI \\ \hline
Recognition Burst & repeated dense detections & No crash, logs controlled \\ \hline
DB Write Stress & frequent visit inserts & No lock failure \\ \hline
\end{tabular}
\end{center}
\normalsize

\textbf{Summary:} Performance is suitable for local institutional deployment.
\end{frame}

% Slide 31: Bug Report and Resolution
\begin{frame}{Bug Report and Resolution}
\footnotesize
\begin{center}
\begin{tabular}{|p{0.9cm}|p{3.2cm}|p{4.1cm}|p{1.1cm}|}
\hline
\textbf{ID} & \textbf{Bug} & \textbf{Resolution} & \textbf{Status} \\ \hline
B1 & Tkinter startup \texttt{TclError} & corrected global Tk font declaration & Fixed \\ \hline
B2 & Known users marked unknown & refined matching threshold + label stabilization & Fixed \\ \hline
B3 & RTSP async\_lock decode issue & switched to go2rtc relay + safer RTSP options & Fixed \\ \hline
B4 & Secondary stream not visible & aligned stream names and app source mapping & Fixed \\ \hline
\end{tabular}
\end{center}
\normalsize
\end{frame}









% Slide 32: Deployment
\begin{frame}{Deployment Overview}
    \textbf{Deployment Mode:} Local Network

    \vspace{0.3cm}
    \textbf{Details:}
    \begin{itemize}
        \item Runs on local machines within the institutional network.
        \item Supports webcam and RTSP network camera streams via \textbf{go2rtc relay}.
        \item No cloud services involved; all data stored locally.
        \item Visit logs, screenshots, and face images stored in local SQLite database.
        \item Multi-camera setup supported using local relay stream endpoints.
    \end{itemize}

    \vspace{0.3cm}
    \textbf{Advantages:}
    \begin{itemize}
        \item Faster response time with local processing.
        \item Data privacy maintained; sensitive information never leaves the network.
    \end{itemize}
\end{frame}
% Slide 33: Individual Contributions (Expanded)
\begin{frame}{Individual Contributions}
    \textbf{Team Members and Contributions:}

    \begin{itemize}
        \item \textbf{Benson Eldho} – System architecture, backend integration, visit logging, module integration, troubleshooting, pipeline orchestration, performance monitoring.
        \item \textbf{Alen Alias} – Performance optimization, documentation, database design, configuration tuning, embedding model management.
        \item \textbf{AV Navneeth Sagar} – Testing and validation, live camera integration, security measures, RTSP stream handling, GUI support.
        \item \textbf{Ashlin P J} – GUI development, report generation, research activities, registration workflows, pose capture guidance, user interface enhancements.
    \end{itemize}

    \vspace{0.3cm}
    \textbf{Image Suggestion:} Optional team photo or individual avatars to visually represent contributions.
\end{frame}
% Slide 34: Demonstration of Working Model (Image Only)
\begin{frame}{Demonstration of Working Model}
    \begin{center}
        \includegraphics[width=\linewidth]{demo.png} 
        % Replace 'demo.png' with your actual screenshot of the working system
    \end{center}
\end{frame}



\begin{frame}{Screenshots of the System}
\begin{center}
    \includegraphics[width=0.8\linewidth]{screenshots.png}
    % Replace 'screenshots.png' with your actual screenshot image
\end{center}
\end{frame}
\begin{frame}{Screenshots of the System}
\begin{center}
    \includegraphics[width=0.8\linewidth]{screenshots.png}
    % Replace 'screenshots.png' with your actual screenshot image
\end{center}
\end{frame}

\begin{frame}{Screenshots of the System}
\begin{center}
    \includegraphics[width=0.8\linewidth]{screenshots.png}
    % Replace 'screenshots.png' with your actual screenshot image
\end{center}
\end{frame}

% Slide 36: Project Budget
\begin{frame}{Project Budget}
    \begin{center}
        \begin{tabular}{|l|c|}
            \hline
            \textbf{Item} & \textbf{Estimated Cost (INR)} \\
            \hline
            Webcam / Camera & 3,000 \\
            \hline
            Computer / GPU Support & 25,000 \\
            \hline
            Software & 0 (Open-source) \\
            \hline
            Storage / Database setup & 1,500 \\
            \hline
            Miscellaneous (Cables, Mounts, etc.) & 1,000 \\
            \hline
            \textbf{Total} & \textbf{30,500} \\
            \hline
        \end{tabular}
    \end{center}

    \vspace{0.3cm}
    \textbf{Note:} Costs are approximate and can be adjusted based on actual purchases.
\end{frame}
% Slide 37: Comparison with Existing Works

\begin{frame}{Comparison with Existing Works}
    \textbf{Proposed System vs Existing Systems:}

    \begin{itemize}
        \item \textbf{Real-time Detection:} Proposed system supports full live detection; many existing systems are partial or offline.
        \item \textbf{Face Recognition:} High accuracy in proposed system; existing works may lack recognition or have lower accuracy.
        \item \textbf{Student Visit Logging:} Only our system maintains timestamped logs for canteen visits.
        \item \textbf{Registration Flexibility:} Supports live capture, image upload, and bulk CSV + photo import.
        \item \textbf{GUI Interface:} User-friendly interface provided; most existing systems lack GUI or have partial support.
        \item \textbf{Database Support:} SQLite database ensures organized storage; existing systems often rely on CSV or local files.
        \item \textbf{Network Stream Support:} RTSP + go2rtc relay support for stable multi-camera access.
        \item \textbf{Performance:} 15-25 FPS; higher than many existing implementations.
    \end{itemize}

    \vspace{0.3cm}
    \textbf{Conclusion:} Our system offers better usability, accuracy, and real-time monitoring specifically for college canteen deployment.
\end{frame}
% Slide 38: Conclusion
\begin{frame}{Conclusion}
    \begin{itemize}
        \item Successfully developed a \textbf{real-time crowd presence and identity detection system} for a college canteen.
        \item Enables \textbf{accurate visit logging} and reduces manual verification errors.
        \item Provides \textbf{GUI and CLI interfaces}, supporting webcam and stable RTSP relay streams.
        \item Employs detection, SFace-based recognition, and multi-sample registration workflows.
        \item Supports bulk onboarding through \textbf{CSV + photo folder import}.
        \item Scalable and \textbf{maintainable architecture} with modular services and database support.
        \item Future scope: integration with \textbf{analytics dashboards, automated reports, and enhanced security measures}.
    \end{itemize}

    \vspace{0.3cm}
    \textbf{Key Takeaway:} The system improves operational efficiency, ensures accountability, and can be adapted to other campus monitoring scenarios.
\end{frame}
% Slide 39: Bibliography
\begin{frame}{Bibliography}
    \begin{itemize}
        \item R. Girshick, \textit{Fast R-CNN}, 2015.
        \item A. Taigman et al., \textit{DeepFace: Closing the Gap to Human-Level Performance in Face Verification}, 2014.
        \item Ultralytics YOLOv8 Documentation, \url{https://docs.ultralytics.com/}
        \item OpenCV-Python Tutorials, \url{https://docs.opencv.org/}
        \item TensorFlow Documentation, \url{https://www.tensorflow.org/}
        \item DeepFace Library Documentation, \url{https://github.com/serengil/deepface}
        \item SQLite Documentation, \url{https://www.sqlite.org/docs.html}
        \item College canteen project references & internal documentation
    \end{itemize}
\end{frame}
\end{document}
